% Generated draft: requires mathematical and source review.
\documentclass[11pt]{article}
\usepackage[margin=25mm]{geometry}
\usepackage{fontspec}
\setmainfont{Latin Modern Roman}
\usepackage{amsmath,amssymb,mathtools,unicode-math}
\setmathfont{Latin Modern Math}
\usepackage{enumitem,needspace,xurl,hyperref}
\setlength{\parindent}{0pt}
\setlength{\parskip}{5pt}
\setlength{\emergencystretch}{3em}
\newcommand{\N}{\mathbb{N}}
\newcommand{\Z}{\mathbb{Z}}
\newcommand{\Q}{\mathbb{Q}}
\newcommand{\R}{\mathbb{R}}
\newcommand{\C}{\mathbb{C}}
\newcommand{\F}{\mathbb{F}}
\newcommand{\A}{\mathbb{A}}
\newcommand{\PP}{\mathbb{P}}
\newcommand{\E}{\mathbb{E}}
\newcommand{\eps}{\varepsilon}
\newcommand{\dd}{\,\mathrm{d}}
\begin{document}
\section*{UnsolvedMath ID: 701164}
\textbf{Source problem code:} catalogue entry
\subsection{Definitions and mathematical statement}
A variety of groups is a class of groups defined by group identities. For a variety \(\Theta\), let \(\Theta^0\) be the category whose objects are the groups free in \(\Theta\) on finitely many generators and whose morphisms are all group homomorphisms between these objects.

An automorphism \(\varphi:\Theta^0\to\Theta^0\) of this category is called \emph{inner} if it is naturally isomorphic to the identity functor. Explicitly, this means that there exists a family of group isomorphisms
\[
 s_A:A\longrightarrow\varphi(A)
 \qquad(A\in\operatorname{Ob}(\Theta^0))
\]
such that, for every pair of objects \(A,B\) and every homomorphism \(\mu:A\to B\),
\[
 \varphi(\mu)=s_B\circ\mu\circ s_A^{-1}.
\]
Write \(P(\Theta)\) for the assertion that every automorphism of \(\Theta^0\) is inner.

For a group \(G\), define its derived series by
\[
 G^{(0)}=G,\qquad G^{(j+1)}=[G^{(j)},G^{(j)}],
\]
where \([H,H]\) is the subgroup generated by the commutators \(x^{-1}y^{-1}xy\) with \(x,y\in H\). For each integer \(d\geq1\), let
\[
 \mathcal S_d=\{G:G^{(d)}=\{1\}\}
\]
be the variety of all groups of derived length at most \(d\). The variety of all metabelian groups is \(\mathcal S_2\); in particular, this includes groups whose derived length is less than two.

The remaining question in source part \(b\) asks whether \(P(\Theta)\) holds for the varieties of solvable groups, and whether it holds for the variety of all metabelian groups. The metabelian question has the precise formulation
\[
 P(\mathcal S_2)\;?
\]
A precise fixed-derived-length interpretation of the solvable question is
\[
 \text{For every integer }d\geq1,\quad P(\mathcal S_d)\;?
\]
This latter formulation is conditional on interpreting the source as referring to the varieties of all groups of bounded derived length. The source does not explicitly specify the derived-length parameter or whether more general varieties consisting of solvable groups are intended. The class of all solvable groups with no uniform bound on derived length is not itself a variety, so it cannot be substituted silently for \(\Theta\).

For context, the supplied source reports that \(P(\Theta)\) holds for the variety of all groups and for the variety of all abelian groups. Archived part \(a\) originally asked whether it holds for the variety of all nilpotent groups of class at most two. The supplied archive reports an affirmative answer, even for nilpotent groups of any fixed class \(n\), attributed to A. Tsurkov. The updated statement of part \(b\) explicitly reports the result for the varieties of all nilpotent groups of class at most \(d\), for \(d\geq2\).

The supplied registry labels the combined record \emph{partially solved}, with the nilpotent part archived as solved and the remaining questions retained in the main section. A separate supplied main-section comment claims an affirmative answer but expressly labels it an \emph{unverified AI claim} and an \emph{unverified AI solution}. That comment is not treated here as an established resolution.
\subsection{Short English statement}
The question asks whether every automorphism of the category of finitely generated free metabelian groups is naturally isomorphic to the identity, and asks the analogous question for varieties of solvable groups. The solvable formulation needs clarification about the intended varieties; one precise interpretation uses each fixed bound on derived length. The source reports the nilpotent part as solved and separately includes an expressly unverified AI claim about an affirmative answer.
\subsection{Sources}
\begin{itemize}
\item Local UnsolvedMath registry, permanent ID \(701164\), exact source problem code \texttt{catalogue entry}, title \emph{Kourovka Notebook Problem 15.76}; supplied complete source record, upstream revision \texttt{d371ee917ac9647fbb06cd6fc529deec6f3dda82}. The supplied status is \texttt{partially\_solved}.
\item B. I. Plotkin, proposer of Problem 15.76, \emph{The Kourovka Notebook (Unsolved Problems in Group Theory)}, 15th issue (2002), with the supplied record describing an update in the 21st issue (2026), arXiv:1401.0300v48. The supplied collection description names E. I. Khukhro and V. D. Mazurov as editors.
\item G. Mashevitzky, B. Plotkin, and E. Plotkin, \emph{J. Algebra} \textbf{282} (2004), 490--512. Cited by the supplied statement for the assertion concerning the variety of all groups; no article title or link is supplied.
\item A. Tsurkov, \emph{Int. J. Algebra Comput.} \textbf{17} (2007), 1273--1281. Cited by the supplied statement and archive for the nilpotent result; no article title or link is supplied.
\item The supplied main-section comment attributes an expressly unverified AI solution to M. Kinyon and J. Urban, \emph{Preprint} (2026), with the supplied link \url{https://kourovkanotebookorg.wordpress.com/wp-content/uploads/2026/09/kourovka-mathematics.pdf}. The attribution and affirmative claim are reproduced as source data, without verification.
\item The supplied status review is dated \texttt{2026-10-16}, records evidence as \texttt{published}, and describes the archived part as solved with remaining parts in the main section. This assessment is supplied dated curation, not independently verified literature assessment. No cited source was visited for this draft.
\end{itemize}
\end{document}
